Considera la funció
\[
g(x)=\begin{cases}
  \dfrac{x^2-4}{x+2} & \si{x<1},\\[8pt]
  x^2+k & \si{x\ge 1},
\end{cases}
\]
on $k$ és un nombre real.

\begin{apartats}

\apartat{0,75}
Calcula $\lim_{x\to-2}g(x)$ i $\lim_{x\to0}g(x)$.

\begin{solucio}
Tots dos punts són a la branca $x<1$.\\
En $x=-2$ hi ha una indeterminació $\tfrac00$:
$\dfrac{x^2-4}{x+2}=\dfrac{(x+2)(x-2)}{x+2}=x-2\to-4$.
El límit val $-4$, tot i que $g(-2)$ no existeix.\\
En $x=0$ se substitueix directament: $\dfrac{-4}{2}=-2$.
\end{solucio}

\apartat{1,25}
Calcula, en funció de $k$, els límits laterals de $g$ en $x=1$. Per a quin valor de $k$
existeix $\lim_{x\to1}g(x)$? Quant val aquest límit?

\begin{solucio}
$\lim_{x\to1^-}g(x)=\dfrac{1-4}{1+2}=-1$ \quad i \quad $\lim_{x\to1^+}g(x)=1+k$.\\
El límit existeix si i només si els dos laterals coincideixen: $1+k=-1$, és a dir
$\boxed{k=-2}$. Aleshores $\lim_{x\to1}g(x)=-1$.
\end{solucio}

\apartat{0,5}
Per a $k=-2$, calcula $\lim_{x\to-\infty}g(x)$ i $\lim_{x\to+\infty}g(x)$.

\begin{solucio}
Quan $x\to-\infty$ actua la primera branca, que per a $x\ne-2$ és $x-2\to-\infty$.\\
Quan $x\to+\infty$ actua la segona: $x^2-2\to+\infty$.
\end{solucio}

\end{apartats}
